\newpage
\phantomsection
\label{prf:every-element-is-one-or-a-successor}

\begin{remark*}[Return]
\hyperref[thm:every-element-is-one-or-a-successor]{Return to Theorem}
\end{remark*}

\begin{theorem*}[Every Element Is Either $1$ or a Successor]
Let $(P,S,1)$ be a Peano system. For every $x \in P$, either $x=1$ or there
exists $u \in P$ such that $S(u)=x$.
\end{theorem*}

\begin{proof}
\textbf{Professional Standard Proof.}~\\
Define a predicate $\varphi$ on $P$ by
\[
\varphi(x)\Longleftrightarrow
\bigl(x=1 \lor \exists u \in P\;(S(u)=x)\bigr).
\]
First, $\varphi(1)$ holds because $1=1$.

Now let $n\in P$ and suppose $\varphi(n)$ holds. Since $n\in P$, the element
$S(n)$ is the successor of an element of $P$, namely of $n$. Hence
\[
\exists u\in P\;(S(u)=S(n)),
\]
by taking $u=n$. Therefore $\varphi(S(n))$ holds.

By the induction principle for a Peano system, $\varphi(x)$ holds for every
$x\in P$. Thus for every $x\in P$, either $x=1$ or there exists $u\in P$ such
that $S(u)=x$.
\end{proof}

\begin{proof}
\textbf{Detailed Learning Proof.}~\\
We want to prove a statement about every element of the Peano system. The
natural tool is induction. The property we want to propagate is:
\[
\varphi(x):\quad x=1 \text{ or } x \text{ is the successor of some element of } P.
\]
In symbols,
\[
\varphi(x)\Longleftrightarrow
\bigl(x=1 \lor \exists u\in P\;(S(u)=x)\bigr).
\]

\textbf{Step 1. Verify the base case.}
For $x=1$, the first alternative in the disjunction is true:
\[
1=1.
\]
Therefore $\varphi(1)$ holds. We do not need to show that $1$ has a
predecessor; the statement allows the special case $x=1$.

\textbf{Step 2. Verify the successor step.}
Let $n\in P$, and assume $\varphi(n)$ holds. We must show $\varphi(S(n))$.
But $S(n)$ is visibly a successor: it is the successor of $n$ itself. Since
$n\in P$, we may choose
\[
u=n.
\]
Then $u\in P$ and
\[
S(u)=S(n).
\]
So the second alternative in the definition of $\varphi(S(n))$ holds:
\[
\exists u\in P\;(S(u)=S(n)).
\]
Hence $\varphi(S(n))$ holds.

\textbf{Step 3. Apply induction.}
We have shown that $\varphi(1)$ holds and that
\[
\forall n\in P\;(\varphi(n)\Longrightarrow \varphi(S(n))).
\]
The induction principle for a Peano system therefore gives
\[
\forall x\in P\;\varphi(x).
\]
Unpacking $\varphi$, this says exactly that for every $x\in P$, either
$x=1$ or there exists $u\in P$ such that $S(u)=x$.
\end{proof}

\begin{remark*}[Proof structure]
Induction is applied to the predicate ``is $1$ or is a successor.'' The base
case uses the first branch of the disjunction. The successor step uses the
second branch directly: $S(n)$ is a successor because it is the successor of
$n$.
\end{remark*}

\begin{remark*}[Dependencies]~\\
\begin{itemize}
  \item \hyperref[def:peano-system]{Peano System}
  \item \hyperref[ax:peano-one-in-n]{Distinguished Element Lies in the Underlying Set}
  \item \hyperref[ax:peano-successor-closure]{Closure Under Successor}
  \item \hyperref[thm:induction-principle-for-peano-system]{Induction Principle for a Peano System}
\end{itemize}
\end{remark*}
